\section{Polynomial stable endomorphisms and twists}
\label{e:section}

The dual cohomology in Proposition~\ref{a:resolution} produces a
degree-three extension class over \(T\).  We first determine its
products, then tensor two copies and use symmetric duality to determine
all negative degrees as well.

\begin{proposition}\label{e:polynomial}
Put \(\mathcal J=T\otimes_C R\), where \(R\) is the resolution in
\eqref{a:simple-resolution}.  There is a triangle in the derived
category of \(T\)-modules
\begin{equation}
s[2]\longrightarrow\mathcal J\longrightarrow s
 \xrightarrow{\tau}s[3],
\label{e:triangle}
\end{equation}
and the connecting class generates the Yoneda algebra:
\begin{equation}
\Ext_T^*(s,s)=k[\tau],\qquad |\tau|=3.
\label{e:one-polynomial}
\end{equation}
The complex \(\mathcal J\) is bounded below, has finite projective
\(T\)-terms, and has bounded homology.
\end{proposition}

\begin{proof}
Tensor the sequence of \((T,C)\)-bimodules
\(0\to DC\to T\to C\to0\) with \(R\).  Projectivity of each
\(R_i\) gives a short exact sequence of complexes of left
\(T\)-modules
\begin{equation}
0\longrightarrow DC\otimes_C R
 \longrightarrow\mathcal J\longrightarrow R\longrightarrow0.
\label{e:coefficient-sequence}
\end{equation}
The two outside complexes have action inflated from \(C\).  For a
finite projective left \(C\)-module \(V\), the natural evaluation map
\[
DC\otimes_C V\longrightarrow D\Hom_C(V,C),\qquad
\phi\otimes v\longmapsto\bigl(f\mapsto\phi(f(v))\bigr)
\]
is an isomorphism, first for a finite free module and then for its
summands.  Applying this termwise identifies the left complex in
\eqref{e:coefficient-sequence} with \(D\Hom_C(R,C)\).  By
\eqref{a:dual-cohomology}, its only homology is \(s\) in homological
degree two, with the asserted \(C\)-module structure.  Good truncation
therefore identifies it with \(s[2]\) in the derived category, also
after inflating the action to \(T\).  Fix this identification.
The quotient complex resolves \(s\), so
\eqref{e:coefficient-sequence} gives \eqref{e:triangle}.
Its middle complex has projective \(T\)-terms, and the triangle gives
its bounded homology.

Because \(\mathcal J\) is bounded below and projective, ordinary
induction--restriction adjunction computes
\[
Q_i:=\Hom_{\mathrm D^b(T)}(\mathcal J,s[i])
 =H^i\Hom_T(\mathcal J,s)
 =H^i\Hom_C(R,s)
 =\begin{cases}k&i=0,\\0&i\ne0.\end{cases}
\]
Apply derived Hom into shifts of \(s\) to \eqref{e:triangle}.
With ordinary negative Ext interpreted as zero, the resulting long
exact sequence contains
\begin{equation}
Q_{i-1}\longrightarrow\Ext_T^{i-3}(s,s)
 \xrightarrow{\alpha\mapsto\alpha[3]\tau}
 \Ext_T^i(s,s)\longrightarrow Q_i.
\label{e:polynomial-sequence}
\end{equation}
At \(i=1\) this is \(k\to0\to\Ext_T^1(s,s)\to0\), so the first
Ext group vanishes.  For every \(i\ge2\), the middle map in
\eqref{e:polynomial-sequence} is an isomorphism.  In particular
\(\Ext_T^2(s,s)=0\), and at \(i=3\) it sends the identity of \(s\)
to the nonzero class \(\tau\).  Induction now shows that \(\tau^m\)
is a basis of \(\Ext_T^{3m}(s,s)\) for every \(m\ge0\), and all
other positive degrees vanish.  Since the isomorphisms are the
displayed Yoneda compositions, this proves the algebra assertion
\eqref{e:one-polynomial}.
\end{proof}

We now tensor two copies of this self-extension algebra. Set
\begin{equation}
X=s\otimes s,\qquad
\cH^a=\sHom_E(X,X[a])\quad(a\in\mathbb Z).
\label{e:stable-groups}
\end{equation}
Let \(\tau_1=\tau\otimes1\) and \(\tau_2=1\otimes\tau\) be the
two external extension classes.  We write \(\mathcal P_m\) for the
space of polynomials of total exponent \(m\) in \(\tau_1,\tau_2\);
their cohomological degree is \(3m\).  This notation distinguishes
these polynomial spaces from the projective resolutions used later.

\begin{proposition}\label{e:all-degrees}
The module \(X\) is nonprojective and
\begin{equation}
\cH^{\ge0}=k[\tau_1,\tau_2],\qquad |\tau_1|=|\tau_2|=3.
\label{e:two-polynomial}
\end{equation}
In every integer degree, the only nonzero groups are
\begin{equation}
\cH^{3m}=\mathcal P_m,\qquad
\cH^{-3m-1}=D\mathcal P_m\qquad(m\ge0).
\label{e:profile}
\end{equation}
Under these identifications, postcomposition by a shifted
\(\tau_i\) between negative groups is the transpose of multiplication
by \(\tau_i\) on the polynomial spaces.  The maps
\(\cH^{-1}\to\cH^2\) given by either class are zero.
\end{proposition}

\begin{proof}
Choose a projective \(T\)-resolution \(Q\to s\).  The total complex
of \(Q\otimes Q\) is a projective \(E\)-resolution of \(X\).  There
are finitely many summands in each total degree, and termwise finite
projectivity gives an isomorphism of cochain complexes
\[
\Hom_E(Q\otimes Q,s\otimes s)
 \cong\Hom_T(Q,s)\otimes\Hom_T(Q,s).
\]
The cohomology of this tensor product is the tensor product of the
cohomologies: over a field each complex splits into its homology and
contractible summands.  Tensoring comparison maps of resolutions
computes the external Yoneda products under this isomorphism.
In particular the maps acting on the first and second tensor factors
commute, since all interchange signs are \(+1\) in characteristic
two.  Proposition~\ref{e:polynomial} therefore gives
\(\Ext_E^*(X,X)=k[\tau_1,\tau_2]\) as algebras.

Its degree-three part is nonzero, so \(X\) is not projective.  Also
\(\End_E(X)=k\), and the identity cannot factor through a projective:
such a factorization would exhibit \(X\) as a projective summand.
Thus \(\sEnd_E(X)=k\).  Positive Ext agrees with positive stable Hom
over the symmetric algebra \(E\), proving \eqref{e:two-polynomial}.
Lemma~\ref{h:duality} gives
\[
\cH^{-a-1}\cong D\cH^a\qquad(a\ge0),
\]
which proves the support and dimensions in \eqref{e:profile}.

To retain the actual multiplication maps, choose a nonzero functional
\(\lambda:\cH^{-1}\to k\).  The composition pairing from
Lemma~\ref{h:duality} identifies \(\cH^{-3m-1}\) with
\(D\mathcal P_m\) by
\begin{equation}
(\beta,p)_m
 =\lambda\bigl(p[-3m-1]\beta\bigr),
\qquad
\beta\in\cH^{-3m-1},\quad p\in\mathcal P_m.
\label{e:pairing}
\end{equation}
It is perfect for every \(m\).  For \(m\ge1\), associativity of
composition gives, for \(p\in\mathcal P_{m-1}\),
\begin{equation}
\bigl(\tau_i[-3m-1]\beta,p\bigr)_{m-1}
 =(\beta,p[3]\tau_i)_m.
\label{e:transpose}
\end{equation}
The right-hand side is pairing with polynomial multiplication by
\(\tau_i\).  This proves the transpose assertion with its precise
shift.  For \(m=0\), the target is \(\cH^2=0\), which proves the
boundary assertion as well.
\end{proof}

One can make the negative action especially concrete.  Let
\(\eta_{a,b}\in\cH^{-3(a+b)-1}\) be dual to
\(\tau_1^a\tau_2^b\).  With multiplication here meaning
postcomposition by the appropriately shifted class,
\begin{equation}
\tau_1\eta_{a,b}=
 \begin{cases}\eta_{a-1,b}&a>0,\\0&a=0,\end{cases}
\qquad
\tau_2\eta_{a,b}=
 \begin{cases}\eta_{a,b-1}&b>0,\\0&b=0.\end{cases}
\label{e:negative-action}
\end{equation}
These formulas include the zero maps out of \(\eta_{0,0}\).
Products of two negative-degree classes vanish by the support
\eqref{e:profile}.

For \(H\in k^\times\), define
\begin{equation}
h_H:T\longrightarrow T,\quad (a,\phi)\longmapsto(a,H\phi),
\qquad \sigma_H=h_H\otimes h_H,
\qquad U_H=E_{\sigma_H}.
\label{e:twist-definition}
\end{equation}
Here \(E_{\sigma_H}\) is the regular left module with right action
\(a\cdot b=a\sigma_H(b)\).  These are automorphisms, and \(U_H\)
is projective on both sides.  The tensor functor \(U_H(-)\) is
restriction by \(\sigma_H^{-1}\), as shown by the natural balanced
isomorphism
\begin{equation}
U_H\otimes_E V\xrightarrow{\ \sim\ }
\Res_{\sigma_H^{-1}}V,\qquad
a\otimes v\longmapsto\sigma_H^{-1}(a)v.
\label{e:inverse-restriction}
\end{equation}
The same convention applies to \(T_{h_H}\otimes_T-\).
Let \(\chi:E\to k\) be the character of \(X\).  Since
\(\chi\sigma_H=\chi\), we fix the resulting identification
\(U_H X=X\).  The exact functor \(U_H(-)\) preserves projectives
and acts on all groups \(\cH^a\), with the shift identifications of
Lemma~\ref{h:evaluation}.

To determine this action in negative degrees, we first choose an
explicit generator in degree \(-1\). Once the action on this generator
is known, the composition pairing determines the action on every
negative group.

\begin{lemma}\label{e:socle}
Put
\[
p_f=f\otimes f,\qquad \zeta=f^*\otimes f^*\in E p_f.
\]
For every \(a,b\in E\),
\(a\zeta b=\chi(a)\chi(b)\zeta\).
The map \(i_X:X\to Ep_f\), \(i_X(1)=\zeta\), is an embedding with
image in \((\rad E)p_f=\Omega X\).  Regarded as a map
\(X\to\Omega X\), it is nonzero stably and spans \(\cH^{-1}\).
\end{lemma}

\begin{proof}
The dual basis vector \(f^*\in DC\) is the character of the
\(f\)-simple on \(C\).  The dual bimodule action in
\eqref{a:trivial-extension} shows that multiplication on either side
by an element of \(C\) scales \(f^*\) by that character; the
square-zero ideal \(DC\) kills it on either side.  Tensoring proves
the asserted action on \(\zeta\).  The projective cover
\(Ep_f\to X\), \(ap_f\mapsto\chi(a)\), has kernel
\((\rad E)p_f\) by \eqref{a:radicals}.  The nonzero vector \(\zeta\)
lies in this kernel, so the displayed map is an embedding into
\(\Omega X\).

We verify nonfactorization through every finite projective.  Write
\(J_E=\rad E\).  If \(b\in E\) satisfies \(J_Eb=0\), symmetry gives
\[
t_E(abj)=t_E(jab)=0\qquad(a\in E,\ j\in J_E),
\]
since \(ja\in J_E\).  Nondegeneracy therefore implies \(bJ_E=0\).
For any map \(X\to E\), every vector in its image satisfies this
condition.  Any left-module map \(E\to J_Ep_f\) has the form
\(b\mapsto br\) with \(r\in J_Ep_f\), so its composite with
\(X\to E\) is zero.  The same conclusion holds for finite sums and
summands of \(E\).  Thus all maps \(X\to\Omega X\) factoring
through a finite projective are zero, whereas \(i_X\ne0\).
Finally \(\cH^{-1}=D\mathcal P_0\) is one-dimensional by
Proposition~\ref{e:all-degrees}.
\end{proof}

\begin{proposition}\label{e:twists}
For every \(H\in k^\times\) and every \(m\ge0\), the action of
\(U_H(-)\) on the groups in \eqref{e:profile} is
\begin{equation}
\left.U_H\right|_{\cH^{3m}}=H^{-m}\,1_{\mathcal P_m},
\qquad
\left.U_H\right|_{\cH^{-3m-1}}=H^{m+2}\,1_{D\mathcal P_m}.
\label{e:twist-weights}
\end{equation}
In particular, the negative action is the stated scalar on the entire
\((m+1)\)-dimensional space.
\end{proposition}

\begin{proof}
First twist \eqref{e:coefficient-sequence} by restriction along
\(h_H^{-1}\).  Because \(h_H\) fixes \(C\), the map
\(h_H\otimes1\) identifies its middle complex with the original
one.  It induces the identity on the quotient \(R\) and multiplication
by \(H\) on \(DC\otimes_C R\).  Write \(\tau_H\) for the twist of
\(\tau\).  Passing to the fixed derived identifications yields a
commutative diagram of triangles
\[
\begin{tikzcd}[column sep=large]
s[2]\arrow[r]\arrow[d,"H"']&
\Res_{h_H^{-1}}\mathcal J\arrow[r]
 \arrow[d,"h_H\otimes1"]&
s\arrow[r,"\tau_H"]\arrow[d,"1"]&
s[3]\arrow[d,"H"]\\
s[2]\arrow[r]&\mathcal J\arrow[r]&s\arrow[r,"\tau"]&s[3].
\end{tikzcd}
\]
Naturality of the connecting morphism gives \(H\tau_H=\tau\).
Hence each external class \(\tau_i\) is sent to
\(H^{-1}\tau_i\), and preservation of products proves the first
formula in \eqref{e:twist-weights}.

For the negative part we first compute the action on \(\cH^{-1}\).
Twist the projective presentation
\[
0\longrightarrow\Omega X\longrightarrow Ep_f
 \longrightarrow X\longrightarrow0.
\]
Under \eqref{e:inverse-restriction}, its comparison with the original
presentation is
\[
\sigma_H:\Res_{\sigma_H^{-1}}(Ep_f)\longrightarrow Ep_f.
\]
It is \(E\)-linear, since
\(\sigma_H(\sigma_H^{-1}(a)b)=a\sigma_H(b)\), and induces the
identity on \(X\).  On the kernels it sends
\(\zeta\) to \(H^2\zeta\).  Thus the generator in
Lemma~\ref{e:socle} is sent to \(H^2\) times itself, proving that
\(U_H\) acts by \(H^2\) on \(\cH^{-1}\).

Finally the exact twist functor preserves shifted composition.
Applying it to the product in \eqref{e:pairing}, and using its actions
on \(\mathcal P_m\) and \(\cH^{-1}\), gives
\[
\bigl(U_H(\beta),H^{-m}p\bigr)_m
 =H^2(\beta,p)_m
\qquad
(\beta\in\cH^{-3m-1},\ p\in\mathcal P_m).
\]
Equivalently,
\((U_H(\beta),p)_m=H^{m+2}(\beta,p)_m\) for every \(p\).
Nondegeneracy of the pairing forces
\(U_H(\beta)=H^{m+2}\beta\) for every \(\beta\), as asserted.
\end{proof}
